\begin{afterword}
\summaryhead

Some people argue that creationism and Intelligent Design are excluded from science a priori, because they invoke the supernatural. The author asks whether this means that creationism could be true and still be barred from science class. The author suggests that the argument lets people avoid saying that religious claims have been tested and found false, and notes that exclusion a priori would justify the creationists' complaint of unfair treatment.

The post adopts Richard Carrier's definition: a supernatural explanation appeals to ontologically basic mental things. It then sets out a dilemma. The author holds that non-reductionism is a confusion. A brain made of quarks cannot concretely envision an irreducible universe, so it can never get Bayesian confirmation for irreducibility. Supernaturalism is non-reductionism about minds.

Still, ID should not be excluded a priori. Every irreducible God has a reducible counterpart, such as an alien with advanced technology, that makes the same predictions, and the predictions can be tested. The Red Sea story fails the archaeological test, and biology shows competing and incompetent design. Religious claims are therefore excluded a posteriori. A corollary says that a designer of artificial general intelligence who speaks of religious beliefs respectfully is no expert on reducing mental things to nonmental ones.

\responsehead

The post's conclusion is sound, and it was not new. Supernatural claims that make predictions can be tested, and the best-known ones have failed. Larry Laudan made a similar objection to the Arkansas creation-science ruling of 1982, and Carrier's article, whose definition the post adopts, opens with the same complaint.

The post's account of the other side is incomplete. It assigns them a motive, avoiding a confrontation with religion, on the strength of ``It seems clear enough.'' It asks what they mean by ``supernatural'' and does not report their answer. Carrier does: in science and law, the word is used for ``the untestable.'' On that meaning the rule excludes untestable claims and says nothing about whether they are true. There is also a legal reason for the rule: the United States Constitution bars teaching religion presented as science, not teaching false science. And the best-known ruling to apply the rule, \textsc{Kitzmiller v. Dover}, also found that ID's attacks on evolution had been refuted, the very verdict the post accuses the rule's users of avoiding.

The author conceded part of the post's middle the same day. Its claim that one ``can never get Bayesian confirmation for a hypothesis of irreducibility'' drew an objection from Benja Fallenstein, and in ``Psychic Powers'' the author answered, ``You're right.'' The conclusion about ID does not depend on that claim. The claim it does depend on, that every supernatural model has a reducible counterpart with the same predictions, is kept in the later post.

The thesis that non-reductionism is a confusion is argued mostly by questions. The post asks what observations a non-reductionist universe would produce. Philosophers of emergence have given an answer: evidence of new forces that act only in certain configurations, as the weak interaction was once new. On that answer, strong emergence, higher-level properties with causal powers of their own, has been judged to lack evidence: a verdict from observation, the kind the post recommends for religion.

The corollary judges unnamed designers of AI to be ignorant or lying, on the basis of one attitude, and gives no case.

In short: a sound case, made before by Laudan and by Carrier, for judging supernatural claims by their predictions, aimed at opponents whose meaning of ``supernatural'' it does not report and to whom it assigns a motive without evidence, with a middle step the author conceded the same day.
\end{afterword}
